%% notes-tex/028-gene-essentiality/sections/8-summary.tex
%% The study's conclusions on one page. The figure is written by
%% experiments/028-gene-essentiality/scripts/summary_figure.py and the enrichment behind
%% panels f and g by misclass_enrichment.py; both read the committed result files.

\section{The results on one page\secstatus{todo}}
\label{sec:summary}

The study had two goals, and this section answers them in that order.

\paragraph{Goal 1: the best held-out score, against Flux Cone Learning and FBA.} On the
195 released held-out genes the best model scores 0.955 and the best single model 0.934,
against 0.742 for Flux Cone Learning and 0.640 for FBA on the same genes
(Figure~\ref{fig:summary}\textbf{a}). The margin of the averaged model over the better
single model was fixed in advance and holds on nine of ten other held-out fifths
(Section~\ref{sec:round75}).

\paragraph{Goal 2: a metabolic module that carries to other datasets.} A record in the
database is one genotype in one environment, so a portable module may read one medium
only, and it may not solve a program inside the training loop at database scale. The
form that meets both is an input and no layer: one parsimonious FBA solve per genotype and
medium, run once offline and cached, read as $\log(1 + |v|)$ by a zero-initialized affine
on the logit, beside that medium's growth ratio and six counts read off the
gene-protein-reaction rules. On the ten repeated splits, three seeds each, that
single-medium form scores $+0.066$ over the same cell with no metabolism, higher on ten of
ten splits and on 29 of 30 split-and-seed pairs, and it matches the four-media form
($+0.002$, Section~\ref{sec:round84-98}). Every form that kept a linear program or a
projection inside the training loop scored lower and less reliably than this one
(Figure~\ref{fig:summary}\textbf{b}), so the layer is retired and the solve is kept as
data. The store for the 13.5 million records of the solid-growth build is 469,474
programs and 71 minutes on 20 cores (Figure~\ref{fig:summary}\textbf{e}).

\tcfigfit[1.0]{figures/summary_2026-09-20-21-42-51.pdf}{\textbf{Gene essentiality on
  yeast-GEM: the best method, the module form that helps, what scales, and the genes
  still missed.} One color means one thing in every panel: gray a comparator with
  nothing learned here, yellow a learned model reading no solve and no graph, orange the
  transformer reading the metabolic solve, red message passing over the gene graphs,
  purple an equal-weight average of the two families; a hatch marks a linear program
  solved inside the training loop.
  \textbf{a}, Every method on the 195 released held-out genes, best first; the rule-count
  rows are round 80. Learned
  models are five-fold ensembles, the random forest is one fit on the split file, and
  intervals are gene bootstraps.
  \textbf{b}, The forms of the metabolic module on one backbone, the pathway multi-hot
  with ProtT5: the bar is the seed mean of the held-out AUROC and each point one seed.
  The gene-rule row adds six redundancy counts to the layer-free form's vector
  (Section~\ref{sec:round77}); the forced growth column ran on ProtT5 alone, twelve
  seeds.
  \textbf{c}, Wall minutes per cell for the forms of \textbf{b}, as run, with up to three
  cells sharing a card.
  \textbf{d}, The same models trained on the 629 released genes and read on the 4,674
  genome-wide genes that carry no GEM reaction, pooled over seeds with a gene bootstrap;
  the dashed line is chance.
  \textbf{e}, What the layer-free solve costs over the 025 solid-growth build,
  logarithmic: the records, those whose deletion touches no GEM gene and read the
  wild-type row, the distinct programs solved once each, and the lethal and uncertified
  ones among them.
  \textbf{f}, What the misranked genes of the released 195 share, per class: the chance
  that a gene carrying the annotation is misranked worse than one without it, 0.5
  (dashed) being no association, for the transformer with the layer-free solve (circles)
  and graph convolution on GO (squares), filled where the Benjamini-Hochberg $q$ within
  the family is under 0.1.
  \textbf{g}, The same on the genome-wide set, the squares being graph convolution on
  ProtT5 over the union graph. Written by
  \texttt{experiments/028-gene-essentiality/scripts/summary\_figure.py}.}{fig:summary}

\begin{description}
  \item[The best method is two models averaged, and the best single model is a tie.]
    Under the fold protocol on the released 195 genes the equal-weight average of graph
    convolution on GO with the transformer reading the solve scores 0.955 with the solve
    vector and the gene-rule counts beside it (Section~\ref{sec:round80}), 0.954 with
    the cached solve head and 0.953 with the layer-free solve, the maximum of thirteen
    post hoc averages. Alone, the rule-count cell scores 0.934, the layer-free form
    0.920 and graph convolution on GO 0.917, with overlapping intervals, against Flux
    Cone Learning's 0.742 and FBA's 0.640 on the same genes
    (Figure~\ref{fig:summary}\textbf{a}). On ten other held-out fifths of the labeled
    genes, with the average fixed in advance, the average beats the better single model
    on nine (mean $+0.023$) and graph convolution on GO beats the transformer on nine
    (Section~\ref{sec:round75}).
  \item[The module form that helps is the one with no layer, and the rules help beside
    it.] On one backbone the cell
    with no metabolism reads 0.858 on six seeds, FBA's growth ratios on the logit 0.890,
    and the deletion's solved flux vector as $\log(1 + |v|)$ on the logit 0.911 on nine
    (sd 0.007, minimum 0.896; the three extra seeds of each are round 81's controls,
    which is why the paired reads below quote the six they share with their partner). Every form that keeps a linear program in the training
    loop reads lower and less reliably: the cached solve with a flux head 0.882 with one
    seed at 0.772, learned ceilings with a flux head 0.845 with two seeds unfit, the
    gated growth column 0.867, and the forced growth column 0.699 on twelve seeds. The
    layer-free form costs 30 minutes per cell with three cells on a card against 83 to
    100 for the forms that solve in the loop
    (Figure~\ref{fig:summary}\textbf{b} and \textbf{c}). On the ten other held-out
    fifths the vector's gain over no metabolism holds on ten of ten ($+0.047$) and its
    gain over the four growth ratios alone does not separate from zero ($+0.011$, sd
    0.023, seven of ten; Section~\ref{sec:round75}). Six counts of
    gene-protein-reaction redundancy beside the vector read 0.918 on the released split
    against 0.908, and 0.862 against 0.843 on the ten fifths, above on six of six seeds
    and eight of ten splits (Section~\ref{sec:round77}). Against the four ratios alone
    the counts are worth $+0.012$ and the solve $+0.024$ on the released split
    (Section~\ref{sec:round82}) and $+0.010$ and $+0.010$ on the ten fifths, so the two
    halves are complementary and neither is reliably the larger
    (Section~\ref{sec:round84-98}).
  \item[Every metabolic column fits inside one experiment.] The four growth ratios are
    four media, so they read four experiments at once; restricted to the one medium the
    solve store is built on, the full recipe reads 0.916 against 0.918 and the solve-only
    recipe 0.912 against 0.911, while the ratio alone falls from 0.890 to 0.877
    (Section~\ref{sec:round83}). The solve supplies from within one environment what the
    four media were carrying, and the ten fifths agree ($+0.002$ for one medium, above
    on six of ten). Without the solve the restriction is not free: one ratio with the
    rule counts reads 0.888 against 0.902 for four, and 0.030 under the full recipe on
    the ten fifths (Section~\ref{sec:round84-98}).
  \item[The layer-free form scales, and it covers only the genes with a reaction.] Over
    the 13.5 million records of the solid-growth build the transform is 469,474
    programs, solved once in 71 minutes on 20 cores and read as a gather in training
    (Figure~\ref{fig:summary}\textbf{e}). Off the metabolic genes every deletion reads
    one constant row, and the transformer with the layer-free solve scores 0.570 on
    the 4,674 genome-wide genes where graph convolution on ProtT5 scores 0.755 and the
    raw degree in the physical graph 0.740 (Figure~\ref{fig:summary}\textbf{d}). A
    model over the whole database therefore needs both inputs: the solve store for the
    genes with a GEM reaction, and the gene graphs as a message-passing input for the
    rest.
  \item[The two families miss different genes, and the annotations say which.] On the
    released 195 the essential genes the transformer ranks low are the ones FBA calls
    viable: the 19 of 31 with a knockout growth ratio of 1.0 have a mean misrank of
    0.134 against 0.025 for the rest (effect 0.91, $q$ 0.003), and under graph
    convolution on GO the same feature has no effect (0.51). No biological-process term
    reaches $q < 0.1$ among the 31 essential genes under either model. The nonessential
    genes both models rank high sit in steroid biosynthesis (effect 0.85 for the
    transformer and 0.96 for the graph model) and, for the transformer, in nucleoside
    triphosphate metabolism (0.88) and monosaccharide biosynthesis (0.83). On the
    genome-wide set the graph model ranks essential genes low when their degree in the
    physical graph is under the class median (effect 0.70) and when they act in
    cytoskeleton organization (0.79) or a microtubule-based process (0.81), and ranks
    nonessential genes high in ncRNA metabolism (0.90), rRNA processing (0.92), RNA
    modification (0.92) and cytoplasmic translation (0.85)
    (Figure~\ref{fig:summary}\textbf{f} and \textbf{g},
    Table~\ref{tab:misclass-enrichment}). Hypothesis, untested: the graph model scores
    a gene by its neighborhood, and the nonessential members of the ribosome biogenesis
    and translation machinery sit in neighborhoods made of essential genes.
\end{description}

%% notes-tex/028-gene-essentiality/tables/misclass_enrichment.tex
%% GENERATED by experiments/028-gene-essentiality/scripts/notes_tex_tables.py
%% SOURCE: results/misclass_enrichment.json
\begin{table}[H]
  \centering
  \scriptsize
  \caption[What the misranked genes share]{What the misranked genes of each held-out set share, per model and class. A gene's misrank is the fraction of the other class on the wrong side of it, averaged over the model's cells; a term is tested within one class by a one-sided Mann-Whitney $U$ of the annotated genes' misrank against the rest, and the effect is $U$ as a probability, the chance an annotated gene is misranked worse than an unannotated one (0.5 is no association), with the Benjamini-Hochberg $q$ within the family in parentheses. Rows per family: the two named features where both sides hold five genes, then the four best specific biological-process terms (at most 5 percent of the class or 12 genes, near-duplicates removed). Written by \texttt{misclass\_enrichment.py}.}
  \label{tab:misclass-enrichment}
  \begin{tabular}{@{}lrrrr@{}}
    \hline
    term & genes & misrank, annotated & misrank, rest & effect ($q$) \\
    \hline
    \multicolumn{5}{@{}l}{\textit{released 195 (metabolic genes); transformer, layer-free solve; essential, ranked low (31 genes, 44 terms, 1 at $q < 0.1$)}} \\
    \quad FBA knockout growth ratio 1.0 & 19 & 0.134 & 0.025 & 0.91 (0.003) \\
    \quad physical degree under the class median & 15 & 0.103 & 0.082 & 0.56 (0.944) \\
    \quad membrane lipid metabolic process & 6 & 0.119 & 0.085 & 0.68 (0.944) \\
    \quad organonitrogen compound biosynthetic process & 8 & 0.116 & 0.083 & 0.65 (0.944) \\
    \quad carbohydrate derivative biosynthetic process & 8 & 0.084 & 0.095 & 0.56 (0.944) \\
    \quad generation of precursor metabolites and energy & 5 & 0.092 & 0.092 & 0.52 (0.944) \\
    \multicolumn{5}{@{}l}{\textit{released 195 (metabolic genes); transformer, layer-free solve; nonessential, ranked high (164 genes, 211 terms, 70 at $q < 0.1$)}} \\
    \quad FBA knockout growth ratio 1.0 & 152 & 0.082 & 0.214 & 0.44 (0.892) \\
    \quad physical degree under the class median & 81 & 0.058 & 0.124 & 0.40 (0.988) \\
    \quad purine ribonucleoside triphosphate metabolic process & 7 & 0.528 & 0.072 & 0.88 (0.006) \\
    \quad steroid biosynthetic process & 5 & 0.324 & 0.084 & 0.85 (0.025) \\
    \quad monosaccharide biosynthetic process & 5 & 0.472 & 0.080 & 0.83 (0.026) \\
    \quad macromolecule modification & 9 & 0.140 & 0.089 & 0.72 (0.043) \\
    \multicolumn{5}{@{}l}{\textit{released 195 (metabolic genes); graph convolution on GO; essential, ranked low (31 genes, 44 terms, 0 at $q < 0.1$)}} \\
    \quad FBA knockout growth ratio 1.0 & 19 & 0.111 & 0.071 & 0.51 (0.990) \\
    \quad physical degree under the class median & 15 & 0.052 & 0.136 & 0.32 (1.000) \\
    \quad monocarboxylic acid metabolic process & 5 & 0.128 & 0.089 & 0.72 (0.758) \\
    \quad small molecule biosynthetic process & 5 & 0.114 & 0.092 & 0.69 (0.758) \\
    \quad generation of precursor metabolites and energy & 5 & 0.164 & 0.082 & 0.68 (0.758) \\
    \quad purine ribonucleotide metabolic process & 7 & 0.091 & 0.097 & 0.60 (0.773) \\
    \multicolumn{5}{@{}l}{\textit{released 195 (metabolic genes); graph convolution on GO; nonessential, ranked high (164 genes, 211 terms, 57 at $q < 0.1$)}} \\
    \quad FBA knockout growth ratio 1.0 & 152 & 0.095 & 0.099 & 0.52 (0.598) \\
    \quad physical degree under the class median & 81 & 0.070 & 0.120 & 0.46 (0.940) \\
    \quad steroid biosynthetic process & 5 & 0.554 & 0.081 & 0.96 (0.004) \\
    \quad lipid biosynthetic process & 12 & 0.329 & 0.077 & 0.83 (0.004) \\
    \quad macromolecule modification & 9 & 0.191 & 0.090 & 0.75 (0.052) \\
    \quad pyruvate metabolic process & 7 & 0.220 & 0.090 & 0.77 (0.059) \\
    \multicolumn{5}{@{}l}{\textit{genome-wide 4,674 (no GEM reaction); transformer, layer-free solve; essential, ranked low (974 genes, 731 terms, 61 at $q < 0.1$)}} \\
    \quad physical degree under the class median & 476 & 0.404 & 0.462 & 0.43 (1.000) \\
    \quad establishment of protein localization to mitochondrion & 25 & 0.651 & 0.428 & 0.74 (0.004) \\
    \quad inner mitochondrial membrane organization & 5 & 0.869 & 0.431 & 0.95 (0.012) \\
    \quad nucleosome organization & 17 & 0.640 & 0.430 & 0.74 (0.016) \\
    \quad localization within membrane & 42 & 0.562 & 0.428 & 0.65 (0.024) \\
    \multicolumn{5}{@{}l}{\textit{genome-wide 4,674 (no GEM reaction); transformer, layer-free solve; nonessential, ranked high (3,700 genes, 1,427 terms, 59 at $q < 0.1$)}} \\
    \quad physical degree under the class median & 1,810 & 0.495 & 0.375 & 0.61 ($<$0.001) \\
    \quad DNA recombination & 113 & 0.592 & 0.429 & 0.65 ($<$0.001) \\
    \quad homologous recombination & 47 & 0.675 & 0.431 & 0.73 ($<$0.001) \\
    \quad RNA modification & 65 & 0.630 & 0.430 & 0.69 ($<$0.001) \\
    \quad DNA repair & 135 & 0.563 & 0.429 & 0.63 ($<$0.001) \\
    \multicolumn{5}{@{}l}{\textit{genome-wide 4,674 (no GEM reaction); graph convolution on ProtT5; essential, ranked low (974 genes, 731 terms, 252 at $q < 0.1$)}} \\
    \quad physical degree under the class median & 476 & 0.337 & 0.158 & 0.70 ($<$0.001) \\
    \quad cytoskeleton organization & 43 & 0.438 & 0.236 & 0.79 ($<$0.001) \\
    \quad microtubule-based process & 36 & 0.459 & 0.237 & 0.81 ($<$0.001) \\
    \quad mitochondrion organization & 37 & 0.379 & 0.240 & 0.74 ($<$0.001) \\
    \quad attachment of spindle microtubules to kinetochore & 20 & 0.488 & 0.240 & 0.81 ($<$0.001) \\
    \multicolumn{5}{@{}l}{\textit{genome-wide 4,674 (no GEM reaction); graph convolution on ProtT5; nonessential, ranked high (3,700 genes, 1,427 terms, 599 at $q < 0.1$)}} \\
    \quad physical degree under the class median & 1,810 & 0.166 & 0.321 & 0.26 (1.000) \\
    \quad ncRNA metabolic process & 173 & 0.635 & 0.226 & 0.90 ($<$0.001) \\
    \quad rRNA processing & 85 & 0.648 & 0.236 & 0.92 ($<$0.001) \\
    \quad cytoplasmic translation & 98 & 0.551 & 0.237 & 0.85 ($<$0.001) \\
    \quad RNA modification & 65 & 0.717 & 0.237 & 0.92 ($<$0.001) \\
    \hline
  \end{tabular}
\end{table}


\subsection{What was run from the closing list, and what is left out\secstatus{todo}}
\label{sec:left-to-run}

Four experiments were listed as able to change a conclusion above. Two are run, one is
left as a design note, and one is out of scope.

\begin{description}
  \item[Repeated outer splits for the finalists: run.] Ten stratified re-splits of the
    1,121 labeled genes at the released proportions, with the finalists and the average
    fixed in advance (Section~\ref{sec:round75}), later extended to the module's forms
    and to nine seeds per split (Section~\ref{sec:round84-98}).
  \item[A rule-aware input for the essential genes the solve calls viable: run.] Six
    counts read off the gene-protein-reaction rules, beside the vector
    (Section~\ref{sec:round77}), under the fold protocol (Section~\ref{sec:round80}),
    alone (Section~\ref{sec:round82}) and column by column
    (Section~\ref{sec:round84-98}). They help on the fifteen genes holding a reaction of
    their own and hurt on the four an alternative gene set covers.
  \item[The gene graphs as an input to the transformer that reads the solve: not run.]
    Round 76 tried the cheapest form, degree and clustering in the physical graph as
    affine columns (Section~\ref{sec:round76}). A faithful version would have to take
    all nine gene graphs, and the transformer already receives those graphs through its
    attention masks and the graph regularization on its attention, so graph-derived node
    inputs would duplicate them; testing the inputs properly means relaxing the masks and
    the regularizer at the same time, which is a change to the backbone and not to the
    module. The averaged model already recovers what the graph family knows.
  \item[The store on a second phenotype: out of scope.] The endpoint here is
    essentiality. Panel \textbf{e} of Figure~\ref{fig:summary} shows the store is
    affordable over the solid-growth build; whether it helps there belongs to that
    experiment.
\end{description}

The construction in which the binding ceilings differ per genotype
(Section~\ref{sec:open}) stays a design question and is not on this list.
